\documentclass[aspectratio=169,xcolor=dvipsnames]{beamer}
\usetheme{SimpleDarkBlue}

\usepackage[T2A]{fontenc}
\usepackage[utf8]{inputenc}
\usepackage[russian]{babel}
\usepackage{amsmath}
\usepackage{amssymb}
\usepackage{booktabs}
\usepackage{array}
\usepackage{graphicx}
\usepackage{adjustbox}

\title{Первичные эксперименты и развитие реализации}
\subtitle{От ванильного RGD к регуляризации и компактной факторизации}
\author{}
\institute{НИУ ВШЭ}
\date{}

\begin{document}

\begin{frame}{Vanilla RGD}
\footnotesize
\renewcommand{\arraystretch}{1.22}
\begin{center}
\begin{tabular}{@{}>{\raggedright\arraybackslash\bfseries}p{0.46\textwidth}p{0.48\textwidth}@{}}
Минимизация функционала: &
$\displaystyle \min_{X\in \mathcal M_r} f(X)$
\\
Функционал: &
$\displaystyle f(X)=\frac12\lVert \mathcal P_{\Omega}(X-Y)\rVert_F^2$
\\
Оператор проекции: &
$\displaystyle (\mathcal P_\Omega(Z))_{ij}=
\begin{cases}
Z_{ij}, & (i,j)\in\Omega,\\
0, & (i,j)\notin\Omega
\end{cases}$
\\
Многообразие матриц фиксированного ранга: &
$\displaystyle \mathcal M_r = \{X\in\mathbb R^{m\times n}: \operatorname{rank}(X)=r\}$
\\
Евклидов градиент: &
$\displaystyle \nabla f(X)=\mathcal P_{\Omega}(X-Y)$
\\
Риманов градиент: &
$\displaystyle \operatorname{grad} f(X)=\operatorname{Proj}_{T_X\mathcal M_r}\bigl(\nabla f(X)\bigr)$
\\
Шаг метода: &
$\displaystyle X_{k+1}=R_{X_k}\bigl(-\eta_k\,\operatorname{grad} f(X_k)\bigr)$
\end{tabular}
\end{center}
\end{frame}

\begin{frame}{Эксперименты с Vanilla RGD}
\small
\begin{columns}[T]
\column{0.48\textwidth}
\begin{block}{Первый эксперимент}
\begin{itemize}
\item матрица $60\times 60$
\item истинный ранг $4$
\item случайные пропуски
\end{itemize}
\end{block}

\begin{block}{Сравнение}
\begin{itemize}
\item ALS
\item Soft-Impute
\item Vanilla RGD
\end{itemize}
\end{block}

\column{0.48\textwidth}
\begin{block}{Проверяем}
\begin{itemize}
\item точность восстановления
\item скорость сходимости
\item устойчивость к шуму
\item ошибку в ранге
\end{itemize}
\end{block}

\begin{block}{Идея}
понять сильные и слабые стороны базовой реализации
\end{block}
\end{columns}
\end{frame}

\begin{frame}{Vanilla RGD: усреднение по методам}
\small
\textbf{Проверяем:}
\begin{itemize}
\item точность восстановления
\item время работы
\item число итераций
\end{itemize}

\begin{center}
\begin{adjustbox}{max totalsize={\textwidth}{0.56\textheight}}
\begin{tabular}{lrrrr}
\toprule
Метод & $\overline{\mathrm{RMSE}}$ & $\overline{\text{Fro}}$ & $\overline{t}$, c & $\overline{N}_{\mathrm{iter}}$ \\
\midrule
ALS & 0.0725 & 0.7287 & 0.0791 & 48.9 \\
Riemannian GD & 0.0454 & 0.4470 & 0.4407 & 195.6 \\
Soft-Impute & 0.0423 & 0.4215 & 0.1460 & 110.2 \\
\bottomrule
\end{tabular}
\end{adjustbox}
\end{center}

\textbf{Вывод:}
\begin{itemize}
\item точность очень хорошая
\item время заметно больше
\item число итераций тоже больше
\end{itemize}
\end{frame}

\begin{frame}{Vanilla RGD: ошибка в ранге}
\scriptsize
\textbf{Проверяем:}
\begin{itemize}
\item что происходит, если ранг задан неправильно
\end{itemize}

\begin{center}
\begin{adjustbox}{max totalsize={\textwidth}{0.58\textheight}}
\begin{tabular}{llrr}
\toprule
case & Метод & $\overline{\mathrm{RMSE}}$ & $\overline{t}$, c \\
\midrule
correct & ALS & 0.0678 & 0.0775 \\
correct & RGD & 0.0413 & 0.4451 \\
correct & Soft-Impute & 0.0426 & 0.1451 \\
\midrule
over & ALS & 0.0723 & 0.1023 \\
over & \textcolor{red}{RGD} & \textcolor{red}{0.0528} & 0.4889 \\
over & Soft-Impute & 0.0428 & 0.1455 \\
\midrule
under & ALS & 0.0775 & 0.0576 \\
under & \textcolor{orange}{RGD} & \textcolor{orange}{0.0422} & 0.3880 \\
under & Soft-Impute & 0.0417 & 0.1474 \\
\bottomrule
\end{tabular}
\end{adjustbox}
\end{center}

\textbf{Вывод:}
\begin{itemize}
\item Vanilla RGD чувствителен к ошибке в ранге
\end{itemize}
\end{frame}

\begin{frame}{Vanilla RGD: шум}
\scriptsize
\textbf{Проверяем:}
\begin{itemize}
\item что происходит при увеличении шума
\end{itemize}

\begin{center}
\begin{adjustbox}{max totalsize={\textwidth}{0.54\textheight}}
\begin{tabular}{llrr}
\toprule
$\sigma$ & Метод & $\overline{\mathrm{RMSE}}$ & $\overline{t}$, c \\
\midrule
0.00 & ALS & 0.0421 & 0.0551 \\
0.00 & RGD & 0.0279 & 0.4212 \\
0.00 & Soft-Impute & 0.0371 & 0.1204 \\
\midrule
0.02 & ALS & 0.0607 & 0.0757 \\
0.02 & RGD & 0.0370 & 0.4305 \\
0.02 & Soft-Impute & 0.0391 & 0.1301 \\
\midrule
0.05 & ALS & 0.1148 & 0.1066 \\
0.05 & \textcolor{red}{RGD} & \textcolor{red}{0.0714} & 0.4702 \\
0.05 & \textcolor{green!50!black}{Soft-Impute} & \textcolor{green!50!black}{0.0508} & 0.1875 \\
\bottomrule
\end{tabular}
\end{adjustbox}
\end{center}

\textbf{Вывод:}
\begin{itemize}
\item Когда матрица шумная или ранг задан неправильно, ограничения fixed-rank manifold мешают хорошо приблизить матрицу.
\end{itemize}

\textbf{Идея:}
\begin{itemize}
\item добавить регуляризацию, как в Soft-Impute
\end{itemize}
\end{frame}

\begin{frame}{Regularized RGD}
\footnotesize
\renewcommand{\arraystretch}{1.22}
\begin{center}
\begin{tabular}{@{}>{\raggedright\arraybackslash\bfseries}p{0.46\textwidth}p{0.48\textwidth}@{}}
Минимизация функционала: &
$\displaystyle \min_{X\in \mathcal M_r} f_{\lambda}(X)$
\\
Функционал: &
$\displaystyle f_{\lambda}(X)=\frac12\lVert \mathcal P_{\Omega}(X-Y)\rVert_F^2+\frac{\lambda}{2}\lVert X\rVert_F^2$
\\
Евклидов градиент: &
$\displaystyle \nabla f_{\lambda}(X)=\mathcal P_{\Omega}(X-Y)+\lambda X$
\\
Риманов градиент: &
$\displaystyle \operatorname{grad} f_{\lambda}(X)=\operatorname{Proj}_{T_X\mathcal M_r}\bigl(\nabla f_{\lambda}(X)\bigr)$
\\
Шаг метода: &
$\displaystyle X_{k+1}=R_{X_k}\bigl(-\eta_k\,\operatorname{grad} f_{\lambda}(X_k)\bigr)$
\\
Параметр регуляризации: &
$\displaystyle \lambda > 0$
\end{tabular}
\end{center}
\end{frame}

\begin{frame}{Regularized RGD: подбор $\lambda$}
\scriptsize
\textbf{Проверяем:}
\begin{itemize}
\item как выбор $\lambda$ влияет на точность, время и число итераций
\end{itemize}

\begin{center}
\begin{adjustbox}{max totalsize={\textwidth}{0.52\textheight}}
\begin{tabular}{rrrr}
\toprule
$\lambda$ & $\overline{\mathrm{RMSE}}$ & $\overline{t}$, c & $\overline{N}_{\mathrm{iter}}$ \\
\midrule
0.0001 & 0.0460 & 0.1292 & 192.2 \\
0.0003 & 0.0457 & 0.1357 & 192.1 \\
0.0010 & 0.0448 & 0.1355 & 191.8 \\
0.0030 & 0.0429 & 0.1298 & 190.7 \\
0.0100 & 0.0397 & 0.1291 & 187.7 \\
\textcolor{green!50!black}{0.0300} & \textcolor{green!50!black}{0.0393} & \textcolor{green!50!black}{0.1172} & \textcolor{green!50!black}{182.1} \\
\bottomrule
\end{tabular}
\end{adjustbox}
\end{center}

\textbf{Вывод:}
\begin{itemize}
\item при усреднении по всем сценариям лучше всего работает $\lambda = 0.03$
\end{itemize}
\end{frame}

\begin{frame}{Regularized RGD: шум и ошибка в ранге}
\scriptsize
\textbf{Проверяем:}
\begin{itemize}
\item как выбор $\lambda$ ведет себя при шуме и ошибке в ранге
\end{itemize}

\begin{columns}[T]
\column{0.48\textwidth}
\centering
\textbf{шум}

\vspace{0.05cm}
\begin{adjustbox}{max totalsize={\linewidth}{0.34\textheight}}
\begin{tabular}{lrrr}
\toprule
$\sigma$ & $\lambda=0.001$ & $\lambda=0.01$ & $\lambda=0.03$ \\
\midrule
0.00 & 0.0265 & \textcolor{green!50!black}{0.0258} & 0.0302 \\
0.02 & 0.0398 & \textcolor{green!50!black}{0.0357} & 0.0361 \\
0.05 & 0.0681 & 0.0578 & \textcolor{green!50!black}{0.0515} \\
\bottomrule
\end{tabular}
\end{adjustbox}

\column{0.48\textwidth}
\centering
\textbf{rank case}

\vspace{0.05cm}
\begin{adjustbox}{max totalsize={\linewidth}{0.34\textheight}}
\begin{tabular}{lrrr}
\toprule
case & $\lambda=0.001$ & $\lambda=0.01$ & $\lambda=0.03$ \\
\midrule
correct & 0.0386 & \textcolor{green!50!black}{0.0339} & 0.0356 \\
over & 0.0577 & 0.0505 & \textcolor{green!50!black}{0.0479} \\
under & 0.0382 & 0.0348 & \textcolor{green!50!black}{0.0343} \\
\bottomrule
\end{tabular}
\end{adjustbox}
\end{columns}

\textbf{Вывод:}
\begin{itemize}
\item при сильном шуме и ошибке в ранге более сильная регуляризация работает устойчивее
\end{itemize}
\end{frame}

\begin{frame}{Regularized RGD: $\lambda = 0.03$}
\scriptsize
\textbf{Проверяем:}
\begin{itemize}
\item как regularized RGD с выбранным $\lambda$ выглядит на фоне других методов
\end{itemize}

\begin{center}
\begin{adjustbox}{max totalsize={\textwidth}{0.50\textheight}}
\begin{tabular}{lrrr}
\toprule
Метод & $\overline{\mathrm{RMSE}}$ & $\overline{t}$, c & $\overline{N}_{\mathrm{iter}}$ \\
\midrule
ALS & 0.0711 & 0.1015 & 50.0 \\
Vanilla RGD & 0.0462 & 0.1559 & 192.3 \\
Soft-Impute & 0.0415 & 0.1110 & 109.9 \\
\textcolor{green!50!black}{RGD + $\lambda = 0.03$} & \textcolor{green!50!black}{0.0393} & 0.1172 & 182.1 \\
\bottomrule
\end{tabular}
\end{adjustbox}
\end{center}

\textbf{Вывод:}
\begin{itemize}
\item регуляризация улучшает точность по сравнению с Vanilla RGD и Soft-Impute
\item дальше нужно решать проблему полной SVD на каждой итерации
\end{itemize}
\end{frame}

\begin{frame}{Compact RGD}
\footnotesize
\begin{center}
\textbf{Bart Vandereycken}

\vspace{0.08cm}
\textit{Low-rank Matrix Completion by Riemannian Optimization}

\vspace{0.08cm}
SIAM Journal on Optimization, 23(2):1214--1236, 2013
\end{center}

\vspace{0.1cm}
\renewcommand{\arraystretch}{1.25}
\begin{center}
\begin{tabular}{@{}>{\raggedright\arraybackslash\bfseries}p{0.34\textwidth}p{0.60\textwidth}@{}}
Факторизация: &
$\displaystyle X = U S V^{\top}$
\\
Касательное направление: &
$\displaystyle \xi = U\, dS\, V^{\top} + U^{\perp} V^{\top} + U (V^{\perp})^{\top}$
\\
Ортогональность: &
$\displaystyle U^{\top}U^{\perp} = 0,\qquad V^{\top}V^{\perp} = 0$
\\
QR-разложения: &
$\displaystyle U^{\perp} = Q_U R_U,\qquad V^{\perp} = Q_V R_V$
\\
Сумма $X + \xi$: &
$\displaystyle
X + \xi =
[U\;\;Q_U]
\begin{bmatrix}
S + dS & R_V^{\top} \\
R_U & 0
\end{bmatrix}
[V\;\;Q_V]^{\top}$
\\
Ретракция: &
$\displaystyle \text{SVD только ядра размера } \le 2r\times 2r$
\end{tabular}
\end{center}

\textbf{Вывод:}
\begin{itemize}
\item шаг становится дешевле, чем в Vanilla RGD
\end{itemize}
\end{frame}

\begin{frame}{Compact RGD: по методам}
\scriptsize
\textbf{Проверяем:}
\begin{itemize}
\item сохранилась ли точность regularized RGD
\item удалось ли ускорить метод
\end{itemize}

\begin{center}
\begin{adjustbox}{max totalsize={\textwidth}{0.50\textheight}}
\begin{tabular}{lrrr}
\toprule
Метод & $\overline{\mathrm{RMSE}}$ & $\overline{t}$, c & $\overline{N}_{\mathrm{iter}}$ \\
\midrule
ALS & 0.0720 & 0.1728 & 49.1 \\
Vanilla RGD & 0.0462 & 0.1559 & 192.3 \\
Soft-Impute & 0.0420 & 0.1211 & 110.3 \\
Compact RGD & 0.0440 & 0.0786 & 194.6 \\
\textcolor{green!50!black}{Compact RGD + reg} & \textcolor{green!50!black}{0.0395} & \textcolor{green!50!black}{0.0697} & \textcolor{green!50!black}{182.0} \\
\bottomrule
\end{tabular}
\end{adjustbox}
\end{center}

\textbf{Вывод:}
\begin{itemize}
\item Compact RGD + reg сохраняет точность regularized RGD
\item время работы заметно меньше
\end{itemize}
\end{frame}

\begin{frame}{Compact RGD: шум и ошибка в ранге}
\scriptsize
\textbf{Проверяем:}
\begin{itemize}
\item остается ли улучшение при шуме и ошибке в ранге
\end{itemize}

\begin{columns}[T]
\column{0.48\textwidth}
\centering
\textbf{шум}

\vspace{0.05cm}
\begin{adjustbox}{max totalsize={\linewidth}{0.28\textheight}}
\begin{tabular}{lrrr}
\toprule
$\sigma$ & Vanilla RGD & Compact RGD & Compact RGD + reg \\
\midrule
0.00 & 0.0270 & 0.0376 & \textcolor{green!50!black}{0.0308} \\
0.02 & 0.0412 & 0.0419 & \textcolor{green!50!black}{0.0365} \\
0.05 & 0.0703 & 0.0626 & \textcolor{green!50!black}{0.0512} \\
\bottomrule
\end{tabular}
\end{adjustbox}

\column{0.48\textwidth}
\centering
\textbf{rank case}

\vspace{0.05cm}
\begin{adjustbox}{max totalsize={\linewidth}{0.28\textheight}}
\begin{tabular}{lrrr}
\toprule
case & Vanilla RGD & Compact RGD & Compact RGD + reg \\
\midrule
correct & 0.0399 & 0.0372 & \textcolor{green!50!black}{0.0358} \\
over & 0.0596 & 0.0539 & \textcolor{green!50!black}{0.0475} \\
under & 0.0391 & 0.0409 & \textcolor{green!50!black}{0.0351} \\
\bottomrule
\end{tabular}
\end{adjustbox}
\end{columns}

\textbf{Вывод:}
\begin{itemize}
\item компактная версия сохраняет выигрыш regularized RGD
\item ускорение достигается за счет SVD ядра размера не больше $2r\times 2r$
\end{itemize}
\end{frame}

\end{document}
